\section{文档目的与架构基线}

本文给出「寻迹」校园失物招领平台的整体架构说明。系统由两个 Vue 前端、
FastAPI 主后端、独立 AI 服务、MySQL、MinIO 与容器化运行环境构成，本文
覆盖系统边界、模块职责、数据组织、核心流程、关键决策与部署方案。架构设计以
校园失物招领业务闭环为中心，同时兼顾可维护性、可靠性、安全性和可演进性。

\begin{table}[H]
  \centering
  \caption{架构说明范围}
  \begin{tabularx}{\textwidth}{L{3.0cm}Y Y}
    \toprule
    说明范围 & 核心内容 & 设计重点 \\
    \midrule
    系统边界 & 用户端、管理端、主后端、AI 与外部服务 & 参与者、容器和通信协议 \\
    数据与状态 & 核心实体、约束、索引和状态机 & 数据一致性与业务闭环 \\
    关键流程 & 发布、匹配、认领、交接与通知 & 事务边界与异步协作 \\
    模块设计 & 分层模块、DTO、持久任务与安全机制 & 低耦合、可测试和可维护 \\
    运行拓扑 & 容器、端口、健康检查和配置 & 可复现部署与服务治理 \\
    \bottomrule
  \end{tabularx}
\end{table}

\subsection{架构目标与驱动因素}

平台面向高校内「信息分散、人工检索效率低、认领核验与隐私保护要求高」的问题，核心闭环为
「发布---检索与匹配---分级验证---双确认交接---归档与信用记录」。在课程项目规模下，
架构优先保证边界清晰、流程可追踪和环境可复现，而不是追求互联网规模。主要架构驱动如下：

\begin{itemize}
  \item 五人团队和课程迭代周期要求降低分布式协调、运维和联调成本；
  \item AI 能力具有不稳定、耗时和依赖较重的特征，不能成为 P0 主流程的单点故障；
  \item 认领、交接、信用和治理操作涉及多实体一致性，需要明确事务与并发守卫；
  \item 证件、凭证和图片需要严格保护，存储、签名、展示与 AI 获取采用分层授权；
  \item 前端、后端、数据和部署之间需要稳定契约，以支持多人并行协作。
\end{itemize}

\section{系统上下文与边界}

\subsection{参与者与外部依赖}

\begin{longtable}{L{2.8cm} L{4.4cm} L{7.0cm}}
  \caption{系统参与者与协作关系}\label{tab:context-actors}\\
  \toprule
  参与者/系统 & 主要目标 & 与寻迹平台的交互 \\
  \midrule
  \endfirsthead
  \toprule
  参与者/系统 & 主要目标 & 与寻迹平台的交互 \\
  \midrule
  \endhead
  普通用户 & 找回失物或归还拾获物 & 注册登录、发布、搜索、查看匹配、发起认领、确认交接、接收通知 \\
  后勤/安保人员 & 批量登记和代管拾获物 & 以 STAFF 身份发布招领、参与审核与交接 \\
  系统管理员 & 治理与运营 & 认证和内容审核、全状态认领审计、举报处置、公告、用户状态、统计和操作日志查询 \\
  MinIO & 保存非结构化资产 & 保存图片和凭证；bucket 保持私有，后端按权限动态签发读取 URL \\
  AI 服务 & 提供辅助计算 & 分类、文本匹配、认领回答语义校验、敏感检测；不直接连接业务数据库 \\
  DashScope/兼容模型服务 & 增强智能推理 & 由 AI 服务统一调用，并配合确定性规则保证稳定输出 \\
  \bottomrule
\end{longtable}

学校统一身份、商业短信、支付、物流和实时聊天位于系统边界之外。验证码认证通过
可替换适配器接入，AI 模型也作为可替换增强能力，使基础发布和认领流程不依赖特定
身份供应商或模型供应商。

\subsection{系统容器图}

图~\ref{fig:container} 展示运行时容器及主要协议。两个前端镜像均采用多阶段构建：Node
构建 Vue 静态资源，Nginx 提供 SPA 回退并将相对路径 \path{/api/} 代理到主后端。

\begin{figure}[H]
  \centering
  \resizebox{0.82\textwidth}{!}{%
  \begin{tikzpicture}[x=1cm,y=1cm]
    \node[actor] (user) at (0,3.5) {普通用户/STAFF\\移动或桌面浏览器};
    \node[actor] (admin) at (0,0.3) {管理员\\桌面浏览器};

    \node[component] (userweb) at (3.2,3.5) {用户端容器\\Vue 3 + Nginx\\HTTPS / api};
    \node[component] (adminweb) at (3.2,0.3) {管理端容器\\Vue 3 + Nginx\\HTTPS / api};
    \node[component] (backend) at (7.2,1.9) {主后端容器\\FastAPI + SQLAlchemy\\业务 API};
    \node[datastore] (mysql) at (11.2,4.2) {MySQL 8.4\\业务数据/任务};
    \node[datastore] (minio) at (11.2,1.9) {MinIO\\私有对象存储};
    \node[component] (ai) at (11.2,-0.4) {AI 服务容器\\智能计算/规则增强};
    \node[component] (model) at (15.2,-0.4) {模型提供方\\DashScope/兼容接口};

    \draw[flow] (user) -- (userweb);
    \draw[flow] (admin) -- (adminweb);
    \draw[flow] (userweb) -- (backend);
    \draw[flow] (adminweb) -- (backend);
    \draw[flow] (backend) -- (mysql);
    \draw[flow] (backend) -- (minio);
    \draw[flow] (backend) -- (ai);
    \draw[dashedflow] (ai) -- (model);

    \begin{scope}[on background layer]
      \node[draw=linegray,dashed,rounded corners=4pt,fit=(userweb)(adminweb)(backend)(mysql)(minio)(ai),inner sep=9pt,
        label={[font=\small\bfseries,color=reportgray]above:寻迹系统部署边界}] {};
    \end{scope}
  \end{tikzpicture}}
  \caption{寻迹系统容器与外部依赖}
  \label{fig:container}
\end{figure}

\section{容器职责与部署单元}

\begin{longtable}{L{2.5cm} L{3.5cm} L{4.9cm} L{3.5cm}}
  \caption{部署容器职责}\label{tab:containers}\\
  \toprule
  容器 & 技术与入口 & 关键职责 & 状态/持久性 \\
  \midrule
  \endfirsthead
  \toprule
  容器 & 技术与入口 & 关键职责 & 状态/持久性 \\
  \midrule
  \endhead
  user-web & Vue 3、Pinia、Vue Router、Element Plus、Nginx & 用户发布、搜索、匹配、认领、通知与个人中心 & 前端无业务真相；token 位于浏览器存储 \\
  admin-web & Vue 3、Pinia、Vue Router、Element Plus、Nginx & 审核、认领审计、举报、公告、用户、统计和日志界面 & 权限最终由后端判断 \\
  backend & Python 3.12、FastAPI、SQLAlchemy async、Alembic & 业务规则、鉴权、事务、对象签名、任务 runner、统一响应 & 业务状态写入 MySQL \\
  ai-service & Python 3.12、FastAPI、httpx、Pillow & 分类、敏感检测、文本匹配、回答语义校验与规则降级 & 无业务数据库、无业务持久化 \\
  mysql & MySQL 8.4 & 业务实体、日志、通知、积分、持久任务 & 命名卷 \term{mysql-data} \\
  minio & 固定日期版本 MinIO 镜像 & 图片和凭证对象；禁止匿名 bucket policy & 命名卷 \term{minio-data} \\
  \bottomrule
\end{longtable}

Compose 将 MySQL 与 MinIO 端口绑定到 \term{127.0.0.1}，对外部署时移除二者的宿主机
端口。后端容器先执行 \term{alembic upgrade head}，再启动 Uvicorn。两个前端通过相对
API 路径复用同源代理，减少浏览器跨域配置。运行检查区分进程存活与依赖就绪，数据库、
对象存储等依赖就绪后才接收业务流量。

\section{主后端模块与分层}

\subsection{模块职责}

\begin{longtable}{L{2.2cm} L{5.2cm} L{7.0cm}}
  \caption{主后端模块划分}\label{tab:backend-modules}\\
  \toprule
  模块 & 核心职责 & 主要协作关系 \\
  \midrule
  \endfirsthead
  \toprule
  模块 & 核心职责 & 主要协作关系 \\
  \midrule
  \endhead
  core & 配置、JWT、异常、AI HTTP 客户端、对象存储、管理员引导 & 为所有业务域提供基础设施 \\
  common/db & 响应包、错误码、分页、校验、时间、Base、AsyncSession、ULID & 无业务规则的共享能力 \\
  user & 注册登录、资料、认证、账号状态 & 认证审批、积分、通知、注销级联 \\
  item & 失物/招领、搜索、图片、验证问题、举报提交 & 发布后写 job；为 match/claim 提供物品视图 \\
  job & 持久任务/outbox、领取、执行、重试、回收 & 驱动分类、敏感检测和匹配 \\
  match & 候选遍历、评分、结果 upsert、匹配通知 & 调用 AI；为 claim 提供关联上下文 \\
  claim & 分级验证、审核、申诉、交接和结案 & 编排 item、match、credit、notification、log \\
  credit & 积分变更、边界裁剪、流水幂等 & 在调用方事务中写用户、流水、通知与日志 \\
  notification & 站内通知、未读数、公告公开查询 & 被匹配、认领、积分、治理模块调用 \\
  admin & 认证/内容审核、举报处置、公告、用户、看板 & 跨业务域编排和操作审计 \\
  operation\_\allowbreak log & 关键动作的审计记录与查询 & 与业务状态在同一会话提交 \\
  \bottomrule
\end{longtable}

\subsection{分层约束与模块协作}

每个业务包采用 \path{router/service/repository/models/schemas/deps}：router 负责 HTTP、
依赖和 Schema；service 负责状态守卫、事务和跨域编排；repository 只负责 SQLAlchemy
查询与写入；ORM 与 Pydantic DTO 分离。router 通过 deps 获得 service，不直接导入
repository 或 DB session，以保持调用方向单一。

跨模块协作由 service 统一编排，user、item 和 match 等领域 service 组合相邻能力；
repository 和 model 只承担明确的数据访问职责。该组织方式保持模块内聚，并为匹配、
认领和交接等跨域事务提供统一编排边界。

\subsection{模块与持久任务关系图}

\begin{figure}[H]
  \centering
  \resizebox{0.98\textwidth}{!}{%
  \begin{tikzpicture}[x=1cm,y=1cm]
    \node[anchor=east,font=\small\bfseries,text=reportgray] at (-1.8,2.6) {同步分层};
    \node[component] (router) at (0,2.6) {API Routers\\HTTP/Schema/权限};
    \node[component] (services) at (4,2.6) {领域 Services\\规则/状态/编排};
    \node[component] (repos) at (8,2.6) {Repositories\\SQLAlchemy/CAS/锁};
    \node[datastore] (db) at (12,2.6) {MySQL\\业务表};

    \node[anchor=east,font=\small\bfseries,text=reportgray] at (-1.8,0) {持久任务};
    \node[component] (item) at (0,0) {Item Service\\发布/编辑};
    \node[datastore] (outbox) at (4,0) {durable\_jobs\\PENDING/RUNNING};
    \node[component] (runner) at (8,0) {Lifespan Runner\\领取/退避/回收};
    \node[component] (ai) at (12,0.75) {AI Service\\计算/规则降级};
    \node[component] (domains) at (12,-0.75) {Match / Claim / Credit\\Notification / Log};

    \draw[flow] (router) -- (services);
    \draw[flow] (services) -- (repos);
    \draw[flow] (repos) -- (db);
    \draw[flow] (item) -- (outbox);
    \draw[flow] (outbox) -- (runner);
    \draw[flow] (runner.east) -- (ai.west);
    \draw[flow] (runner.east) -- (domains.west);
  \end{tikzpicture}}
  \parbox{0.96\textwidth}{\zihao{-5}\color{reportgray}
    \textbf{注：}上行表示请求经 Router、Service 和 Repository 访问 MySQL；
    下行表示 Item Service 在业务事务中同步写入 \texttt{durable\_jobs}，Runner
    随后领取或回收任务，调用 AI 计算（失败时规则降级）并执行领域操作。
    结果、通知和日志仍通过上行分层路径持久化。}
  \caption{分层调用与持久任务协作关系}
  \label{fig:module-job}
\end{figure}

\section{数据架构}

\subsection{聚合与实体关系}

核心实体关系以用户、物品、匹配和认领为主轴：用户可发布多条 lost/found；一条 found
可包含多张图片和最多三个验证问题；lost 与 found 的候选关系由唯一 match pair 表示；
claim 可选关联 match，且一条 claim 最多一个 handover；通知、积分、举报和操作日志记录
横切副作用。所有业务主键为 26 位 ULID，避免依赖单库自增并保持大体时间有序。

\begin{table}[H]
  \centering
  \caption{核心数据表与一致性作用}
  \begin{tabularx}{\textwidth}{L{3.5cm}Y Y}
    \toprule
    表组 & 代表表 & 关键约束/用途 \\
    \midrule
    用户与认证 & \texttt{users}、\texttt{user\_cert\_requests} & 手机号唯一；用户状态、角色、认证和信用快照 \\
    物品与资产 & \texttt{lost\_items}、\texttt{found\_items}、\texttt{item\_images}、\texttt{verify\_questions} & 状态、审核、敏感标记；文件仅存稳定引用 \\
    匹配 & \texttt{match\_results} & \texttt{lost\_item\_id + found\_item\_id} 唯一；保存分项分、来源与降级标志 \\
    认领与交接 & \texttt{claim\_requests}、\texttt{claim\_answers}、\texttt{handover\_records} & handover 的 claim 唯一；回答保留问题与参考答案快照，claim 保存校验来源和降级标志 \\
    信用与通知 & \texttt{credit\_logs}、\texttt{notifications} & 积分按用户、业务、原因唯一；通知带关联类型与 ID \\
    治理与审计 & \texttt{reports}、\texttt{announcements}、\texttt{operation\_logs} & 同一举报人对同一精确目标唯一；关键状态动作可追溯 \\
    可靠任务 & \texttt{durable\_jobs} & 类型、业务、版本唯一；状态/领取与业务版本索引 \\
    \bottomrule
  \end{tabularx}
\end{table}

\subsection{索引、约束与迁移}

数据库结构以 Alembic 迁移链作为唯一版本来源。迁移既负责创建或调整表结构，也负责在
增加唯一约束前收敛历史数据；部署过程只执行迁移，不并行维护手写建表脚本。业务关系
采用逻辑 ID 和应用层守卫，由 service 维护引用完整性；唯一性、任务类型和任务状态等
关键不变量同时使用数据库约束保护。迁移前需要核对旧库结构，避免跳过数据收敛步骤。

\section{关键时序：发布、任务与匹配}

图~\ref{fig:publish-sequence} 描述可靠发布路径。接口成功只表示物品与任务已经提交，
不表示 AI 计算已经结束；分类失败时任务重试，AI 不可用时匹配切换为规则评分，
敏感检测未得出明确安全结论时保持保护状态。

\begin{figure}[H]
  \centering
  \resizebox{\textwidth}{!}{%
  \begin{tikzpicture}[x=1cm,y=1cm]
    \node[state,text width=2.7cm] (submit) at (0,2.6) {1. 提交发布\\字段 + assetRef};
    \node[state,text width=2.7cm] (validate) at (4,2.6) {2. 后端校验\\身份/字段/资产};
    \node[state,text width=2.7cm] (write) at (8,2.6) {3. 同一事务\\物品、图片\\日志、jobs};
    \node[state,text width=2.7cm] (return) at (12,2.6) {4. COMMIT\\返回发布成功};

    \node[state,text width=2.7cm] (claimjob) at (12,0) {5. Runner 领取\\锁定持久任务};
    \node[state,text width=2.7cm] (score) at (8,0) {6. 智能计算\\AI 结果\\规则增强};
    \node[state,text width=2.7cm] (persist) at (4,0) {7. 原子写回\\匹配与通知\\任务完成};

    \draw[flow] (submit) -- (validate);
    \draw[flow] (validate) -- (write);
    \draw[flow] (write) -- (return);
    \draw[flow] (return) -- (claimjob);
    \draw[flow] (claimjob) -- (score);
    \draw[flow] (score) -- (persist);
  \end{tikzpicture}}
  \parbox{0.94\textwidth}{\zihao{-5}\color{reportgray}
    Runner 使用行锁领取任务；AI 响应不可用时切换为规则评分，执行异常时按退避策略重试。}
  \caption{物品发布与持久任务执行流程}
  \label{fig:publish-sequence}
\end{figure}

\subsection{匹配算法与可解释性}

综合评分的目标权重为图像 0.4、文本 0.3、地点 0.2、时间 0.1，但只在「双方均有输入且该
维度存在可用评分器」的维度上归一。主后端校验远端返回的分值、来源和可用性，并自行
重新计算总分；图像评分器不可用时通过 \texttt{imageAvailable=false} 排除该维度。规则
评分采用中英文词元 Jaccard、地点相等/包含关系和失物时间区间距离。总分低于 70 不保留
活动匹配，达到 90 且候选为证件时发送高优先级通知。数据库保存 \texttt{degraded} 与
\texttt{scoreSource}，用于区分模型增强结果与规则结果。

\section{事务、并发与状态一致性}

\subsection{事务所有权}

repository 只执行查询、flush 或条件更新，不自行 commit；顶层 service 复用一个
\term{AsyncSession} 编排多表写入，并通过显式事务边界统一提交或回滚。发布流程在一次
提交中写物品、图片、操作日志和 job；任务执行在一次提交中写匹配、通知和任务状态；
交接结案在一次提交中写 claim、found、关联 lost、handover、积分、通知和日志。

\subsection{并发控制手段}

\begin{table}[H]
  \centering
  \caption{并发与幂等控制}
  \begin{tabularx}{\textwidth}{L{3.3cm}Y Y}
    \toprule
    场景 & 设计机制 & 目的 \\
    \midrule
    同一招领被并发认领 & 锁定 found/claim/match；found 使用期望状态 CAS：PENDING 到 CLAIMING & 只有一个申请占用物品 \\
    并发审核 & 行锁加源状态条件更新 & 状态变化后返回冲突，不覆盖新结果 \\
    双方并发确认交接 & 锁定 claim、found、handover；确认位更新后重新读取 & 只在双方均确认时结案一次 \\
    重复交接记录 & handover 的 claim 唯一键，加 savepoint 捕获并发冲突 & 重复创建返回已有记录 \\
    重复积分 & 业务唯一键、用户行锁、savepoint、边界裁剪 & 状态与积分保持一致且可重试 \\
    重复匹配任务 & match pair 唯一、任务业务版本唯一、upsert & 不重复创建 pair 与通知 \\
    多 runner 领取任务 & \term{FOR UPDATE SKIP LOCKED}、RUNNING 租约、版本检查 & 并行领取且进程崩溃后可回收 \\
    \bottomrule
  \end{tabularx}
\end{table}

核心状态使用统一枚举字典：lost 从 SEARCHING 到 FOUND 或 CLOSED；
found 从 PENDING 到 CLAIMING，再到 RETURNED 或 CLOSED；claim 包含 PENDING、ANSWER\_PASSED、
PROOF\_PENDING、APPROVED、REJECTED、APPEALING、HANDED\_OVER 和 TERMINATED。状态变更由
service 校验角色、所有权和源状态，前端只负责展示允许的操作和业务结果。

\section{AI 隔离、降级与可靠性}

AI 服务是独立进程，通过 \path{/internal/ai/*} 提供分类、敏感检测、匹配评分和认领回答语义校验。主后端
使用带 3 秒默认超时和 \term{X-Service-Token} 的 \term{httpx.AsyncClient}；响应必须满足
键集合、类型、有限数值、取值范围和来源白名单，否则按失败处理。AI 服务自身可调用
DashScope 原生或兼容接口，并复用进程级连接池。

\begin{longtable}{L{3.1cm} L{5.6cm} L{6.0cm}}
  \caption{AI 能力与降级策略}\label{tab:ai-fallback}\\
  \toprule
  能力 & 正常路径 & 降级/失败路径 \\
  \midrule
  \endfirsthead
  \toprule
  能力 & 正常路径 & 降级/失败路径 \\
  \midrule
  \endhead
  物品分类 & VL 模型返回五类之一，关键词补充标签 & 关键词分类，标记 \texttt{degraded=true}；无可用结果的持久任务重试 \\
  文本匹配 & 聊天评分或 embedding 余弦相似度 & 中文二元/英文词元 Jaccard，加地点与时间规则 \\
  认领回答 & 小模型计算回答与多个参考答案的语义相似度，批量返回逐题分数 & 关键词命中规则；业务记录保存来源与降级标志，普通用户不获得参考答案 \\
  图像相似度 & 可选视觉评分器返回相似度 & 评分器不可用时排除该维度，不以「有图」代替相似度 \\
  敏感检测 & VL 返回唯一类别或 NONE & 任意超时、空结果、降级、需复核均保持敏感；任务退避重试 \\
  \bottomrule
\end{longtable}

这种策略体现「业务可用性优先、隐私安全从严」：AI 失败不回滚已经提交的基础发布，
也不能将未证明安全的图片公开。敏感检测只负责判定，不生成脱敏副本；无权用户不获得
原图地址，界面使用受控占位内容表示图片不可见。

\section{安全与隐私架构}

\subsection{身份、权限与配置安全}

用户 JWT 解析后仍会查询数据库，只有存在且状态为 ACTIVE 的用户可继续访问，因而旧
token 不能绕过账号禁用；管理员接口使用后端角色依赖，不依赖前端路由守卫。密码采用
PBKDF2-SHA256、随机盐和恒时比较。非 DEBUG 环境若仍使用仓库默认 JWT 或管理员密码，
后端拒绝启动；CORS 不允许通配符，只接受显式 origin 列表。AI 内部接口除健康检查外
均要求至少 32 字符服务令牌，使用恒时比较；空令牌只在明确的本地开发模式允许。
认领参考答案快照不进入普通用户响应，只由管理员治理接口返回；管理端路由守卫仅改善
交互，最终仍由后端 ADMIN 角色依赖授权。

\subsection{图片与对象存储安全}

上传服务不信任 MIME、文件名或扩展名，而是使用 Pillow 解码、限制 10 MiB 和像素数、
拒绝 decompression bomb，再按真实 JPEG/PNG/WebP 格式重编码，以去除 EXIF 和文本
元数据。对象 key 包含业务类型、上传者、年月和随机 UUID；提交业务时只接受当前用户、
对应业务类型的 \term{asset://} 引用，并通过 MinIO stat 验证对象存在。

MinIO bucket 在启动和上传时确保存在并删除匿名 policy。普通与敏感资源使用不同签名
时长；AI 收到的是由内部 MinIO endpoint 独立生成的短签名，不是浏览器预览 URL。AI
下载器还执行协议/凭据/主机白名单检查、DNS 解析后的私网检查、每次重定向复验、大小
限制、内容类型与实际解码格式比对，防范 SSRF 与恶意图片请求。

\section{部署与运行架构}

\subsection{环境划分}

\begin{table}[H]
  \centering
  \caption{环境与部署策略}
  \begin{tabularx}{\textwidth}{L{2.8cm}Y Y}
    \toprule
    环境 & 推荐形态 & 关键约束 \\
    \midrule
    本地开发 & MySQL/MinIO 容器，前后端源码热更新；也可完整 Compose & 允许受控 debug SMS 和 mock/规则 AI，不提交真实秘密 \\
    小组联调 & 固定主机和 API，应用容器化，数据库与对象存储地址稳定 & 固定枚举和字段；使用独立测试账号与数据 \\
    课程演示 & 使用固定联调环境，冻结代码、迁移和演示数据 & 答辩前不临时改表；保留规则降级和固定脚本 \\
    对外部署 & HTTPS 网关，仅公开前端/443，内部化 backend/AI/MinIO & 数据与对象服务不直接暴露到公网 \\
    \bottomrule
  \end{tabularx}
\end{table}

配置通过 Pydantic Settings 和环境变量注入；示例文件进入仓库，真实 \path{.env} 被忽略。
Compose 对数据库、JWT、管理员和 MinIO 密码使用必填展开，AI token 在共享环境应替换为
唯一随机值。镜像安装均使用锁文件和 frozen 模式，前端、后端、AI 均有独立 Dockerfile。

\section{质量属性分析}

\begin{longtable}{L{2.4cm} L{6.1cm} L{6.1cm}}
  \caption{质量属性与架构响应}\label{tab:quality-attributes}\\
  \toprule
    属性 & 设计措施 & 质量响应 \\
  \midrule
  \endfirsthead
  \toprule
    属性 & 设计措施 & 质量响应 \\
  \midrule
  \endhead
  可维护性 & 业务域模块化、router/service/repository、DTO/ORM 分离、文档化枚举和状态机 & 模块职责清晰，接口与数据结构可独立演进 \\
  可靠性 & DB outbox、任务租约、退避、版本去重、AI 规则降级、fail-closed & 异步任务可恢复，外部模型波动不阻塞主流程 \\
  一致性 & 单会话统一提交、行锁、CAS、唯一约束、savepoint & 多实体状态在统一事务中闭环，重复请求保持幂等 \\
  安全性 & JWT 状态复核、角色依赖、私有存储、短签名、SSRF 防护、默认密钥拒启 & 身份、对象和服务调用均采用分层授权 \\
  性能 & 全异步 HTTP/DB、连接池、受控 AI 并发、索引与分页 & 发布请求快速返回，耗时计算由持久任务异步处理 \\
  可部署性 & 锁定依赖、Alembic、Compose、健康检查、多阶段镜像 & 开发、联调和课程演示环境可重复构建 \\
  可审计性 & 统一响应、业务错误码、操作日志、积分流水、匹配来源 & 关键业务动作与状态变化均可追踪 \\
  \bottomrule
\end{longtable}

\section{关键架构决策记录}

\begin{longtable}{L{1.4cm} L{3.5cm} L{5.4cm} L{4.1cm}}
  \caption{架构决策摘要}\label{tab:adr}\\
  \toprule
  ADR & 决策 & 原因与收益 & 约束与边界 \\
  \midrule
  \endfirsthead
  \toprule
  ADR & 决策 & 原因与收益 & 约束与边界 \\
  \midrule
  \endhead
  001 & 主后端采用模块化单体，AI 独立部署 & 控制五人团队的联调、配置和排错成本，同时隔离模型依赖 & 强一致业务模块共享数据库；AI 只通过内部 API 协作 \\
  002 & 主后端与 AI 统一 Python async 栈 & FastAPI/Pydantic 契约清晰，httpx 与异步 DB 适合外部调用 & CPU 密集或阻塞任务不得占用事件循环 \\
  003 & MySQL + 私有 MinIO，库中只存稳定引用 & 结构化数据与大对象分离，可动态授权并缩短敏感 URL & 业务提交校验资产归属；对象读取必须经过授权 \\
  004 & 数据库 outbox 替代仅内存 BackgroundTasks & 发布与任务原子落库，进程重启可恢复；无需引入消息中间件 & runner 只通过持久任务接口领取、重试和提交结果 \\
  005 & AI 可降级，敏感判断 fail-closed & 模型不可用不阻断主流程，又不以不确定结果换取公开展示 & 匹配和认领问答记录来源；高敏感或不确定结果进入人工复核 \\
  006 & 双 Vue SPA 共享 TypeScript 类型 & 用户端移动优先、管理端桌面任务密集，复用枚举和契约 & 共享包只包含跨端类型和常量，不承载页面业务 \\
  007 & ULID 与统一枚举 & 跨模块 ID 一致、可排序；避免中文状态和魔法值扩散 & 接口、数据表与前端类型使用同一状态语义 \\
  \bottomrule
\end{longtable}

\section{结论}

寻迹架构以模块化单体承载强一致业务，以独立 AI 服务承载可替换计算，以 MySQL outbox
保证异步任务可恢复，并通过私有对象存储、动态签名和 fail-closed 保护图片资产。
该架构遵循「发布不被 AI 阻塞、状态变化可追踪、多实体写入统一提交、敏感结果分层展示」
四项核心原则，为校园失物招领从发布到交接归档提供清晰的系统边界和关键设计约束。
